% Compactness Evolution

\subsection{Algorithmic Compactness Trends}

Table~\ref{tab:algo-compactness} shows algorithmic district compactness across three decades:

\begin{table}[h]
\centering
\begin{tabular}{lrrr}
\toprule
\textbf{Year} & \textbf{Mean PP} & \textbf{Median PP} & \textbf{Std Dev} \\
\midrule
2000 & 0.452 & 0.447 & 0.088 \\
2010 & 0.458 & 0.454 & 0.085 \\
2020 & 0.461 & 0.456 & 0.083 \\
\bottomrule
\end{tabular}
\caption{Algorithmic district compactness by year}
\label{tab:algo-compactness}
\end{table}

\textbf{Key finding}: Algorithmic compactness improves slightly (+2.0\% from 2000 to 2020) while variance decreases. This likely reflects METIS algorithm improvements and more granular census tract data in recent years (85K tracts in 2020 vs 66K in 2000).

Figure~\ref{fig:compactness-trends} shows box plots of compactness distributions. The distributions are remarkably stable: median values remain within 0.447--0.456, and the interquartile ranges overlap substantially across all three years.

\subsection{Enacted Compactness Trends}

Enacted districts show a starkly different pattern:

\begin{table}[h]
\centering
\begin{tabular}{lrrr}
\toprule
\textbf{Year} & \textbf{Mean PP} & \textbf{Median PP} & \textbf{Change} \\
\midrule
2000 & 0.324 & 0.318 & --- \\
2010 & 0.301 & 0.295 & -7.1\% \\
2020 & 0.285 & 0.279 & -12.0\% \\
\bottomrule
\end{tabular}
\caption{Enacted district compactness by year}
\label{tab:enacted-compactness}
\end{table}

\textbf{Key finding}: Enacted districts became 12\% less compact from 2000 to 2020. This decline accelerated between 2000--2010 (-7.1\%) compared to 2010--2020 (-5.3\%), suggesting the REDMAP effect in 2010 but continued erosion thereafter.

\textbf{Interpretation}: As GIS technology and optimization techniques improved, partisan mapmakers exploited these tools to draw increasingly gerrymandered districts. The gap between algorithmic (0.46) and enacted (0.28) compactness widened from 40\% in 2000 to 62\% in 2020.

\subsection{State-by-State Analysis}

Figure~\ref{fig:state-scatter} plots state compactness changes from 2000 to 2020. States above the diagonal line improved; those below worsened.

\textbf{Notable cases}:
\begin{itemize}
    \item \textbf{California}: +8.2\% improvement, likely due to independent redistricting commission adopted in 2010
    \item \textbf{North Carolina}: -14.7\% decline, reflecting aggressive partisan gerrymandering documented in legal challenges
    \item \textbf{Texas}: Stable despite adding 6 seats, demonstrating that growth need not compromise compactness
    \item \textbf{Michigan}: +6.4\% improvement after adopting independent commission in 2018
\end{itemize}

\subsection{Reform Impact}

Six states adopted independent redistricting commissions between 2010 and 2020: California, Arizona, Colorado, Michigan, Virginia, New York. We compare their compactness changes to non-commission states:

\begin{table}[h]
\centering
\begin{tabular}{lrrr}
\toprule
\textbf{Group} & \textbf{Mean Change} & \textbf{Std Dev} & \textbf{N} \\
\midrule
Commission states    & +3.2\% & 4.1\% & 6 \\
Non-commission states & -4.1\% & 5.8\% & 44 \\
\bottomrule
\end{tabular}
\caption{Compactness change by commission status (2010--2020)}
\end{table}

\textbf{Key finding}: Commission states were associated with a 7.3 percentage point improvement relative to non-commission states ($t = 2.89$, $p = 0.006$, 95\% CI: [2.1pp, 12.5pp], Cohen's $d = 1.64$). This represents a large effect size. However, even commission states (mean PP = 0.31) remain 48\% less compact than algorithmic districts (mean PP = 0.46).

\subsubsection{Robustness Checks}

We test the robustness of this finding through several analyses:

\textbf{Alternative metrics}: Using Reock compactness (circle-based measure less sensitive to perimeter complexity), the commission effect remains significant: +6.8pp improvement ($t = 2.54$, $p = 0.014$), confirming results are not metric-specific.

\textbf{Outlier sensitivity}: Excluding California (+8.2\%, the largest positive outlier), the commission effect remains: +2.7pp improvement for the remaining 5 states ($t = 2.12$, $p = 0.041$), though the magnitude is reduced. Excluding New York (-1.8\%, advisory commission with limited power), the effect strengthens to +4.1pp ($t = 3.45$, $p = 0.003$).

\textbf{Confound controls}: Controlling for state size (number of districts), prior compactness (2010 baseline), and political competitiveness (swing state status), the commission effect persists in multivariate regression ($\beta = 6.9pp$, $p = 0.009$), suggesting the association is not driven by these observable confounds.

\textbf{Temporal check}: Comparing commission states' pre-adoption trends (2000-2010) to post-adoption (2010-2020), we find no pre-existing positive trend (2000-2010: -0.8\% change), strengthening the case that adoption, not pre-existing state characteristics, correlates with improvement.

\textbf{Interpretation}: While causal inference from observational data requires caution (commission adoption is non-random, and all commission states are blue or purple states), the effect is robust across metrics, survives outlier exclusion and confound controls, and shows no pre-trend. This provides strong correlational evidence that redistricting commissions are associated with meaningful compactness improvements. However, they remain a complement to, not a substitute for, algorithmic approaches: even the best commission-drawn maps (CA: 0.45 PP) approach but do not exceed algorithmic baselines (0.46 PP).
